\section{Introduction}


%Presents a clear learning context and includes:
%1) the general scientific concept with the necessary background;
%2) primary goals and the specific scientific concept to be investigated;
%2) specific theory, experimental design and equations needed with fully defined quantities;
%3) any assumptions made or required;
%4) hypothesis to be investigated;
%5) references to supporting sources.

Electromotive Response, or EMF, is responsible for the mechanical-electrical energy transfer across generators,
motors, and transformers.
The objective of this study is to demonstrate this exchange in the form of magnetic dipoles accelerating
conductive coils.

\begin{equation}\label{eq:faraday}
    \epsilon_1 = -\frac{\dif \phi}{\dif t}
\end{equation}

In Eq.~\eqref{eq:faraday}, Faraday's Law describes that a change in flux results in an EMF\@.

By forming the differential equation that relates flux to the magnetic field, we can show that,

\begin{equation}\label{eq:flux}
    d \phi = \frac{\mu_0 p_m}{2 \pi} \frac{z^2 - r^2 / 2}{(r^2 + z^2)^{5/2}} (2 \pi r dr)
\end{equation}

where $d \phi$ is the flux element, $p_m$ is the magnetic dipole moment, $z$ is the axial distance away from the
coil, $\mu_0$ is the permeability of free space, and $r$ is the radius of the loop.

Further, we can derive the EMF by integrating across the time element.

\begin{equation}\label{eq:emf}
    \epsilon = -\frac{3}{2} N \mu_0 p_m r_0^2 \frac{z}{(r_0^2 + z^2)^{5/2}}v
\end{equation}

Where $\epsilon$ is the EMF, $N$ is the number of coil turns, $r_0^2$ is the coil radius, and $v$ is the speed of
the magnet moving through the coil.

Additionally, for sufficiently large axial distances, Eq. ~\eqref{eq:Bz} applies.
This allows us to determine the magnetic dipole moment from $B_z$ and distance measurements.

\begin{equation}\label{eq:Bz}
    B_z \approx p_m \frac{\mu_0}{2 \pi} \frac{1}{z^3}
\end{equation}


%Ref equation: Eqs.~\eqref{eq:norm_damp}, ~\eqref{eq:crit_damp}, and \eqref{eq:osc_damp}, respectively.